% DIAPOSITIVA 3: El problema
\begin{frame}{El problema: saber quién va a ganar es costoso}
  \begin{columns}[T]
    \begin{column}{0.52\textwidth}
      \textbf{La industria de pronósticos electorales}
      \begin{itemize}
        \item Industria global de opinión pública:\\
              \num{\$153 mil millones} facturados en 2024
        \item Colombia: más de \num{\$700 mil millones COP} (2024)
        \item Plataformas predictivas alternativas:\\
              Polymarket mueve más de \num{1 billón USD} en contratos de predicción
      \end{itemize}
      \vspace{0.5em}
      \textbf{Tres retos del sector en Colombia}
      \begin{itemize}
        \item \alrt{Regulatorio:} Ley 2494/2025 endureció divulgación.
              \textbf{GAD3} se retiró del país en 2026
        \item \alrt{Cobertura:} más de \num{1.100 municipios}; encuestadoras
              cubren solo 10-12 ciudades
        \item \alrt{Costo:} una encuesta multimuestra es inviable
              para elecciones subnacionales
      \end{itemize}
    \end{column}
    \begin{column}{0.44\textwidth}
      \vspace{0.3em}
      \begin{block}{La pregunta natural}
        ¿La actividad ciudadana en redes sociales dice algo sobre cómo va a votar el colombiano?
      \end{block}
      \vspace{0.5em}
      \begin{block}{Definición: \textit{Nowcasting}}
        Estimación del valor \textbf{actual} de una variable subyacente difícil de observar directamente, a partir de señales disponibles en tiempo real.
      \end{block}
    \end{column}
  \end{columns}
\end{frame}

% DIAPOSITIVA 4: Literatura
\begin{frame}{De los modelos econométricos a las redes sociales}
  \begin{itemize}
    \item \textbf{1970s-2000s:} primeros modelos de nowcasting electoral con variables
          económicas y de aprobación (VAR, ARCH/GARCH): alta precisión histórica
    \item \textbf{2016:} golpe de gracia. Ningún modelo predijo \textbf{Trump} ni el \textbf{Brexit}
    \item \textbf{2009:} Tumasjan \textit{et al.}: menciones en Twitter como estimador insesgado
          de participación electoral; MAE inferior al 2\% con millones de tweets
    \item \textbf{2010s-2020s:} modelos híbridos (redes más variables económicas) mejoran resultados
          pero requieren infraestructura de extracción masiva, inviable fuera de academia
    \item \textbf{Brito \& Andeato (2020s):} \kw{SOMEN-DC}: \textit{Social Media Electoral Nowcasting
          During Campaign}. Promete superar encuestadoras con interacciones de candidatos
  \end{itemize}

  \vspace{0.4em}
  \begin{alertblock}{Problema identificado en SOMEN-DC}
    Tres limitaciones estructurales que los autores no resuelven:
    \begin{itemize}
      \item \alrt{Techo encuestocéntrico:} la variable objetivo \textbf{son las encuestas}, no los resultados
      \item \alrt{During-Campaign:} requiere scraping diario desde el primer día de campaña
      \item \alrt{No transferibilidad:} pesos específicos a cada elección; imposible generalizar
    \end{itemize}
  \end{alertblock}
\end{frame}
